% TOP_PROBLEM: {"id":30007528,"problem_number":"LOCAL-30007528","title":"Optimal inflation target","category_id":21,"category":{"id":21,"name":"economics","display_name":"Economics","description":"Open economic theory statements, published research agendas, and proposed scoped specifications; question provenance and historical priority are qualified per record.","slug":"economics","order_index":21,"created_at":"2026-10-08T00:00:00Z"},"status":"provenance_pending","external_url":null,"legacy_ids":[],"published":false,"statement":"Find an empirically robust set of optimal constant inflation targets in a specified heterogeneous monetary economy with wage rigidity and a lower rate bound.","economics":{"original_id":17,"field":"Inflation and monetary policy","kind":"design","formalization_basis":"proposed_specification","source_status":"not_verified","priority_status":"not_established","priority_note":"No explicit primary open-problem statement matching this catalogue item was verified; supporting research is not evidence of first formulation.","earliest_verified_reference":null,"current_reference":{"authors":["Hess Chung","Callum Jones","Antoine Lepetit","Fernando M. Martin"],"title":"Implications of Inflation Dynamics for Monetary Policy Strategies","year":2025,"url":"https://www.federalreserve.gov/econres/feds/files/2025072pap.pdf","doi":"10.17016/FEDS.2025.072","locator":"Abstract, p. 1; Section 1, pp. 2–3","evidence_summary":"The paper studies robust monetary strategies and supply-shock tradeoffs and gives model-based guidance. It does not state the precise minimax persistence-uncertainty problem formulated below as an unresolved conjecture."},"formalization_note":"New model-based welfare specification. The cited source studies inflation strategies, not an explicit unresolved optimal-target conjecture matching this formulation."},"statusQualification":"Provenance pending: an explicit published statement of this exact open question was not verified. This is a catalogue-authored candidate specification, not a certified open conjecture. New model-based welfare specification. The cited source studies inflation strategies, not an explicit unresolved optimal-target conjecture matching this formulation.","statusReviewedAt":"2026-10-08"}
% FIRST_POSTED: 2026-10-08
% ENTRY_KIND: statement_only
% !TeX program = lualatex
\documentclass[11pt]{article}
\usepackage{fontspec}
\usepackage[a4paper,margin=25mm]{geometry}
\usepackage{amsmath,amssymb,mathtools}
\usepackage{xurl}
\usepackage[hidelinks]{hyperref}
\newcommand{\sref}[2]{\hyperlink{source:#1:#2}{[S#2]}}
\setlength{\emergencystretch}{3em}
\begin{document}
% problemId:30007528
\section*{Optimal inflation target}\label{30007528}
\subsection{Definitions and mathematical statement}
Fix a fully specified stochastic monetary economy with household types \(i\), social weights \(\omega_i\), sticky prices, downward nominal wage constraints, and an effective lower bound \(i_t\ge\underline i\) on nominal policy rates. Preferences \(u_i(c,\ell)\), shock laws, fiscal financing, and the admissible class \(\mathcal R(\bar\pi)\) of policy rules around a constant inflation target \(\bar\pi\) are given. Restrict targets to a compact interval \([\pi_L,\pi_U]\) and require a declared equilibrium selection and integrable discounted utility for every evaluated rule. Here \(\beta\in(0,1)\), and nonnegative welfare weights sum to one; assume a maximizing target exists or report only the supremum.
Define
\[
W(\bar\pi)=\sup_{r\in\mathcal R(\bar\pi)}
E_r\sum_{t\ge0}\beta^t\sum_i\omega_i u_i(c_{it},\ell_{it}),\qquad
\bar\pi^*\in\arg\max_{\bar\pi\in[\pi_L,\pi_U]}W(\bar\pi).
\]
Determine an empirically credible set for \(\bar\pi^*\) when wage rigidity, neutral-rate dynamics, and household inflation exposure are only partially identified. Quantify the tradeoff between lower-bound stabilization benefits, price dispersion, and heterogeneous purchasing-power costs within the same welfare criterion. Report whether the ranking of targets is robust to admissible parameter uncertainty and policy-rule restrictions. An optimum calculated in one calibrated model is a conditional answer, not a universal resolution. This specification asks for a defensible robust target in a declared environment; social weights and institutional policy choices cannot be inferred solely from inflation data.

\par\textbf{Formulation and scope.} New model-based welfare specification. The cited source studies inflation strategies, not an explicit unresolved optimal-target conjecture matching this formulation.

\par\textbf{Open-question provenance.} An explicit published open-question statement was not verified. Sources below are supporting literature and must not be treated as a first listing.

\par\textbf{Historical priority.} No explicit primary open-problem statement matching this catalogue item was verified; supporting research is not evidence of first formulation.

\subsection{Short English statement}
Find an empirically robust set of optimal constant inflation targets in a specified heterogeneous monetary economy with wage rigidity and a lower rate bound.

\subsection{Sources}
\begin{itemize}\raggedright
\item[\textnormal{[S1]}]\hypertarget{source:30007528:1}{} Hess Chung, Callum Jones, Antoine Lepetit, Fernando M. Martin. 2025. Implications of Inflation Dynamics for Monetary Policy Strategies. Abstract, p. 1; Section 1, pp. 2–3.
\url{https://www.federalreserve.gov/econres/feds/files/2025072pap.pdf}.
DOI: 10.17016/FEDS.2025.072.
{\small\textit{Supporting or current literature; see evidence qualification. The paper studies robust monetary strategies and supply-shock tradeoffs and gives model-based guidance. It does not state the precise minimax persistence-uncertainty problem formulated below as an unresolved conjecture.}}
\end{itemize}
\end{document}
